\section{Quasi-subsequential finite-state transducers}

\begin{definition}
  A \emph{$k$-separated quasi-subsequential finite-state transducer} $\mathcal{T}$ is
  a tuple
  \begin{equation*}
    \langle \Sigma, \Omega, Q, X, s, F, \tau, \rho, \eta, \eta_x, \lambda, \psi
    \rangle
  \end{equation*}
  for $k > 0$ where
  \begin{itemize}
    \item $Q, s, F, \tau, \rho, \eta, \lambda$ and $\psi$ are tuples of $k$ elements,
    \item $\mathcal{T}_i = \langle \Sigma, \Omega, Q_i, s_i, F_i, \tau_i, \lambda_i,
	    \psi_i \rangle$ is a subsequential finite-state transducer for $i \in \{
	    1, 2, \ldots, k \}$,
    \item $X$ is a separator set such that $X \subseteq Q_i$ for $i \in \{ 1, 2,
            \ldots, k \}$,
    \item $\rho_i : Q_i \rightarrow X$ is a quasi-state reduction function such that
	    $dom(\rho_i) = rng(\tau_i)$ for $i \in \{ 1, 2, \ldots, k \}$,
    \item $\eta_i : Q_i \times X \rightarrow Q_i$ is a quasi-state extension function
	    for $i \in \{ 1, 2, \ldots, k \}$,
    \item $\eta_x : X \times X \rightarrow X$ is a separator extension function.
  \end{itemize}
\end{definition}

\begin{definition}
  A state in a $k$-separated quasi-subsequential finite-state transducer $\mathcal{T}
  = \langle \Sigma, \Omega, Q, X, s, F, \tau, \rho, \eta, \eta_x, \lambda, \psi
  \rangle$ is an element of the set $\prod_{i=1}^k Q_i$.
\end{definition}

We write $q^i$ for the $i$-th element of a state $q \in \prod_{i=1}^k Q_i$.

\begin{definition}
  Let $\mathcal{T} = \langle \Sigma, \Omega, Q, X, s, F, \tau, \rho, \eta, \eta_x,
  \lambda, \psi \rangle$ be a $k$-separated quasi-subsequential finite-state
  transducer. The \emph{iterated separator extension functions} $\eta_x^i : X^i
  \rightarrow X$ are defined as
  \begin{itemize}
    \item $\eta_x^1(x) = x$,
    \item $\eta_x^i(x_1, x_2, \ldots, x_i) = \eta_x(\eta_x^{i-1}(x_1, x_2, \ldots,
	    x_{i-1}), x_i)$ for $i \in \{ 2, 3, \ldots, k \}$.
  \end{itemize}
\end{definition}

\begin{definition}
  Let $\mathcal{T} = \langle \Sigma, \Omega, Q, X, s, F, \tau, \rho, \eta, \eta_x,
  \lambda, \psi \rangle$ be a $k$-separated quasi-subsequential finite-state
  transducer. The \emph{separator transition function} $\delta_x : \prod_{i=1}^k Q_i
  \times \Sigma \rightarrow X$ is defined as
  \begin{equation*}
    \delta_x(q, \sigma) = \eta_x^k(\rho_1(\tau_1(q^1, \sigma)), \rho_2(\tau_2(q^2,
    \sigma)), \ldots, \rho_k(\tau_k(q^k, \sigma))).
  \end{equation*}
\end{definition}

\begin{definition}
  Let $\mathcal{T} = \langle \Sigma, \Omega, Q, X, s, F, \tau, \rho, \eta, \eta_x,
  \lambda, \psi \rangle$ be a $k$-separated quasi-subsequential finite-state
  transducer. The \emph{quasi-state transition function} $\delta_i : \prod_{i=1}^k Q_i
  \times \Sigma \rightarrow Q_i$ is defined as
  \begin{equation*}
    \delta_i(q, \sigma) = \eta_i(\tau_i(q^i, \sigma), \delta_x(q, \sigma))
  \end{equation*}
  for $i \in \{ 1, 2, \ldots, k \}$.
\end{definition}

\begin{definition}
  Let $\mathcal{T} = \langle \Sigma, \Omega, Q, X, s, F, \tau, \rho, \eta, \eta_x,
  \lambda, \psi \rangle$ be a $k$-separated quasi-subsequential finite-state
  transducer. The \emph{state transition function} $\delta : \prod_{i=1}^k Q_i \times
  \Sigma \rightarrow \prod_{i=1}^k Q_i$ is defined as
  \begin{equation*}
    \delta(q, \sigma) = \langle \delta_1(q, \sigma), \delta_2(q, \sigma), \ldots,
    \delta_k(q, \sigma) \rangle.
  \end{equation*}
\end{definition}

\begin{definition}
  Let $\mathcal{T} = \langle \Sigma, \Omega, Q, X, s, F, \tau, \rho, \eta, \eta_x,
  \lambda, \psi \rangle$ be a $k$-separated quasi-subsequential finite-state
  transducer. The \emph{generalized state transition function} $\delta^* :
  \prod_{i=1}^k Q_i \times \Sigma^* \rightarrow \prod_{i=1}^k Q_i$ is defined as
  \begin{itemize}
    \item $\delta^*(q, \epsilon) = q$ for $q \in \prod_{i=1}^k Q_i$,
    \item $\delta^*(q, w\sigma) = \delta(\delta^*(q, w), \sigma)$ for $q \in
	    \prod_{i=1}^k Q_i$, $w \in \Sigma^*$ and $\sigma \in \Sigma$.
  \end{itemize}
\end{definition}

\begin{definition}
  Let $\mathcal{T} = \langle \Sigma, \Omega, Q, X, s, F, \tau, \rho, \eta, \eta_x,
  \lambda, \psi \rangle$ be a $k$-separated quasi-subsequential finite-state
  transducer. The \emph{state transition output function} $\lambda : \prod_{i=1}^k Q_i
  \times \Sigma \rightarrow \Omega^*$ is defined as
  \begin{equation*}
    \lambda(q, \sigma) = \bigwedge_{i=1}^k{\lambda_i(q^i, \sigma)}.
  \end{equation*}
\end{definition}

\begin{definition}
  Let $\mathcal{T} = \langle \Sigma, \Omega, Q, X, s, F, \tau, \rho, \eta, \eta_x,
  \lambda, \psi \rangle$ be a $k$-separated quasi-subsequential finite-state
  transducer. The \emph{generalized state transition output function} $\lambda^* :
  \prod_{i=1}^k Q_i \times \Sigma^* \rightarrow \Omega^*$ is defined as
  \begin{itemize}
    \item $\lambda^*(q, \epsilon) = \epsilon$ for $q \in \prod_{i=1}^k Q_i$,
    \item $\lambda^*(q, w\sigma) = \lambda^*(q, w)\lambda(\delta^*(q, w), \sigma)$
	    for $q \in \prod_{i=1}^k Q_i$, $w \in \Sigma^*$ and $\sigma \in \Sigma$.
  \end{itemize}
\end{definition}

\begin{proposition}
  Let $\mathcal{T} = \langle \Sigma, \Omega, Q, X, s, F, \tau, \rho, \eta, \eta_x,
  \lambda, \psi \rangle$ be a $k$-separated quasi-subsequential finite-state
  transducer. Then
  \begin{equation*}
    \lambda(\delta^*(q, u), \sigma)^{-1} \lambda^*(q, u)^{-1} u = \lambda^*(q, u
    \sigma)^{-1} v
  \end{equation*}
  for $q \in \prod_{i=1}^k Q_i$, $u \in \Sigma^*$, $\sigma \in \Sigma$ and $v \in
  \Omega^*$.
\end{proposition}

\begin{proof}
  \begin{align*}
    \lambda(\delta^*(q, u), \sigma)^{-1} \lambda^*(q, u)^{-1} v
     & = (\lambda^*(q, u) \lambda(\delta^*(q, u), \sigma))^{-1} v \\
     & = \lambda^*(q, u \sigma)^{-1} v.
  \end{align*}
\end{proof}

\begin{definition}
  A \emph{path in a $k$-separated quasi-subsequential finite-state transducer}
  $\mathcal{T} = \langle \Sigma, \Omega, Q, X, s, F, \tau, \rho, \eta, \eta_x,
  \lambda, \psi \rangle$ is either trivial or non-trivial. A trivial path $\pi_q$ in
  $\mathcal{T}$ is associated with every state $q \in \prod_{i=1}^k Q_i$. The length of
  $\pi_q$ is $0$. The start state of $\pi_q$ is $q$. The end state of $\pi_q$ is $q$.
  The label of $\pi_q$ is $\langle \epsilon, \epsilon \rangle$. The path $\pi_q$ is
  succsessful if and only if $q = s$ and $q^i \in F_i$ for some $i \in \{ 1, 2,
  \ldots, k \}$. A non-trivial path $\pi$ in $\mathcal{T}$ is a finite sequence of
  state transitions
  \begin{align*}
    \langle q_1, \sigma_1, w_1, q_2 \rangle, \langle q_2, \sigma_2, w_2, q_3 \rangle,
    \dots, \langle q_{n-1}, \sigma_{n-1}, w_{n-1}, q_n \rangle
  \end{align*}
  for $n > 1$ where $q_{i+1} = \delta(q_i, \sigma_i)$ and $w_i = \lambda(q_i,
  \sigma_i)$ for $i \in \{ 1, 2, \ldots, n-1 \}$. The length of $\pi$ is $n-1$. The
  start state of $\pi$ is $q_1$. The end state of $\pi$ is $q_n$. The label of $\pi$
  is $\langle \sigma_1 \sigma_2 \ldots \sigma_{n-1}, w_1 w_2 \ldots w_{n-1} \rangle$.
  The path $\pi$ is successful if and only if $q_1 = s$ and $q_n^i \in F_i$ for some
  $i \in \{ 1, 2, \ldots, k \}$.
\end{definition}

\begin{proposition}
  Let $\mathcal{T} = \langle \Sigma, \Omega, Q, X, s, F, \tau, \rho, \eta, \eta_x,
  \lambda, \psi \rangle$ be a $k$-separa\-ted quasi-subsequential finite-state
  transducer. There exists a path in $\mathcal{T}$ with start state $q' \in
  \prod_{i=1}^k Q_i$, end state $q'' \in \prod_{i=1}^k Q_i$ and label $\langle u, v \rangle
  \in \Sigma^* \times \Omega^*$ if and only if $\delta^*(q', u) = q''$ and
  $\lambda^*(q', u) = v$.
\end{proposition}

\begin{proof}
  We prove the statement by induction on the length of $u$. If $|u| = 0$, then $u =
  \epsilon$. Let there be a path with start state $q' \in \prod_{i=1}^k Q_i$, end state
  $q'' \in \prod_{i=1}^k Q_i$ and label $\langle u, v \rangle = \langle \epsilon, v
  \rangle$. Only trivial paths have a component of their label equal to $\epsilon$.
  Therefore, $q' = q''$, $v = \epsilon$ and by definition $\delta^*(q', u) =
  \delta^*(q', \epsilon) = q' = q''$ and $\lambda^*(q', u) = \lambda^*(q', \epsilon)
  = \epsilon = v$. A similar argument follows in the other direction. Otherwise, if
  $|u| > 0$, then $u = u'\sigma$ for $u' \in \Sigma^*$ and $\sigma \in \Sigma$. Let
  there be a path with start state $q' \in \prod_{i=1}^k Q_i$, end state $q''' \in
  \prod_{i=1}^k Q_i$ and label $\langle u, v \rangle$. There exist a state $q'' \in
  \prod_{i=1}^k Q_i$ and a path with start state $q'$, end state $q''$ and label
  $\langle u', v' \rangle$ where $v' \in pref(v)$. By definition $\delta(q'', \sigma)
  = q'''$ and $\lambda(q'', \sigma) = v'^{-1}v$. By the inductive hypothesis
  $\delta^*(q', u') = q''$ and $\lambda^*(q', u') = v'$. Therefore, by definition
  $\delta^*(q', u) = \delta^*(q', u'\sigma) = \delta(\delta^*(q', u'), \sigma) =
  \delta(q'', \sigma) = q'''$ and $\lambda^*(q', u) = \lambda^*(q', u'\sigma) =
  \lambda^*(q', u')\lambda(\delta^*(q', u'), \sigma) = v' \cdot \lambda(q'', \sigma)
  = v' \cdot v'^{-1}v = v$. Again, a similar argument can be followed in the other
  direction.
\end{proof}

\begin{definition}
  The \emph{function recognized by a $k$-separated quasi-subsequential finite-state
  transducer} $\mathcal{T} = \langle \Sigma, \Omega, Q, X, s, F, \tau, \rho, \eta,
  \eta_x, \lambda, \psi \rangle$ is the set of all labels of successful paths
  \begin{align*}
    \mathcal{F}_\mathcal{T} = \{ \langle w, \lambda^*(s, w) \psi_i(\delta^*(s, w)^i)
    \rangle \mid \exists i \in \{ 1, 2, \ldots, k \}(\delta^*(s, w)^i \in F_i) \}.
  \end{align*}
\end{definition}

\iffalse
  \begin{theorem}
    Let $\mathcal{T}$ be a trimmed functional $k$-separated finite-state
    transducer such that $\mathcal{F}_\mathcal{T}$ has bounded variation. Let
    $\mathcal{T'}$ be the quasi-subsequen\-tial $k$-separated finite-state
    transducer derived from $\mathcal{T}$ by applying the construction described
    above. Then $\mathcal{F}_\mathcal{T'} = \mathcal{F}_\mathcal{T}$.
  \end{theorem}

  \begin{proof}  
    We define $p_{h+1} = \delta(p_h, \sigma_{i_h}) =
      \langle p_{h+1}^1, p_{h+1}^2, \ldots, p_{h+1}^k \rangle$ where
    \begin{align*}
      p_{h+1}^j & = \delta_j(p_h, \sigma_{i_h}) \\
                & = \begin{aligned}[t]
         & \begin{aligned}[t] \{
           & \langle p''', \lambda(p_h,
          \sigma_{i_h})^{-1}uvw \rangle \mid                      \\
           & \exists \langle p', u \rangle \in p_h^j              \\
           & \exists \langle p', \sigma_{i_h}, v, p'' \rangle \in
          \Delta                                                  \\
           & \exists \langle p''', w \rangle \in
          \Delta_\epsilon(p'')(p''' \in P_j) \} \cup
        \end{aligned} \\
         & \bigcup_{r = 1}^k{
        \begin{aligned}[t] \{
           & \langle p''', \lambda(p_h,
          \sigma_{i_h})^{-1}uvw \rangle \mid          \\
           & \exists \langle p', u \rangle \in p_h^r  \\
           & \exists \langle p', \sigma_{i_h}, v, p''
          \rangle \in \Delta                          \\
           & \exists \langle p''', w \rangle \in
          \Delta_\epsilon(p'')(p''' \in P_r) \}
        \end{aligned}
        }
      \end{aligned}  \\
                & = \bigcup_{r = 1}^k{
      \begin{aligned}[t] \{
         & \langle p''', \lambda(p_h,
        \sigma_{i_h})^{-1}uvw \rangle \mid          \\
         & \exists \langle p', u \rangle \in p_h^r  \\
         & \exists \langle p', \sigma_{i_h}, v, p''
        \rangle \in \Delta                          \\
         & \exists \langle p''', w \rangle \in
        \Delta_\epsilon(p'')(p''' \in P_j) \}
      \end{aligned}
      }
    \end{align*}
    for $j \in \interval{1}{k}$.
    If $h < l$ then $\sigma_{i_{h+1}} \in \Sigma$ and if
    $i_{h+1} - i_h > 1$ then $\sigma_j = \epsilon$ for $j \in
      \interval{i_h+1}{i_{h+1}-1}$. Therefore,
    \begin{align*}
      \langle q_{i_{h+1}}, w_{i_h} w_{i_h+1} \ldots w_{i_{h+1}-1} \rangle
      \in \{
      \langle p''', vw \rangle \mid
       & \exists \langle q_{i_h}, \sigma_{i_h}, v, p'' \rangle \in
      \Delta                                                          \\
       & \exists \langle p''', w \rangle \in \Delta_\epsilon(p'') \}.
    \end{align*}
    By the inductive hypothesis we have
    \begin{gather*}
      \langle q_{i_{h+1}}, \Delta_\wedge(\sigma_{i_1}
      \sigma_{i_2} \ldots \sigma_{i_{h-1}})^{-1} w_1 w_2 \ldots
      w_{i_{h+1}-1} \rangle \in\\
      \bigcup_{r=1}^k \begin{aligned}[t] \{
        \langle p''', uvw \rangle \mid
         & \exists \langle p', u \rangle \in p_h^r                      \\
         & \exists \langle p', \sigma_{i_h}, v, p'' \rangle \in \Delta  \\
         & \exists \langle p''', w \rangle \in \Delta_\epsilon(p'') \}.
      \end{aligned}
    \end{gather*}
    Substituting the left-hand side further we get
    \begin{gather*}
      \langle q_{i_{h+1}}, \lambda(p_1, \sigma_{i_1} \sigma_{i_2} \ldots
      \sigma_{i_{h-1}})^{-1} w_1 w_2 \ldots w_{i_{h+1}-1} \rangle \in\\
      \bigcup_{r=1}^k \begin{aligned}[t] \{
        \langle p''', uvw \rangle \mid
         & \exists \langle p', u \rangle \in p_h^r                      \\
         & \exists \langle p', \sigma_{i_h}, v, p'' \rangle \in
        \Delta                                                          \\
         & \exists \langle p''', w \rangle \in \Delta_\epsilon(p'') \}.
      \end{aligned}
    \end{gather*}
    Removing the common output prefix we get
    \begin{gather*}
      \langle q_{i_{h+1}}, \lambda(p_h, \sigma_{i_h})^{-1} \lambda(p_1,
      \sigma_{i_1} \sigma_{i_2} \ldots \sigma_{i_{h-1}})^{-1} w_1 w_2
      \ldots w_{i_{h+1}-1} \rangle \in\\
      \bigcup_{r=1}^k \begin{aligned}[t] \{
        \langle p''', \lambda(p_h, \sigma_{i_h})^{-1} uvw
        \rangle \mid
         & \exists \langle p', u \rangle \in p_h^r                      \\
         & \exists \langle p', \sigma_{i_h}, v, p'' \rangle \in
        \Delta                                                          \\
         & \exists \langle p''', w \rangle \in \Delta_\epsilon(p'') \}.
      \end{aligned}
    \end{gather*}
    By \Cref{prop:lambda_derivative} we get
    \begin{gather*}
      \langle q_{i_{h+1}}, \lambda(p_1, \sigma_{i_1} \sigma_{i_2} \ldots
      \sigma_{i_h})^{-1} w_1 w_2 \ldots w_{i_{h+1}-1} \rangle \in\\
      \bigcup_{r=1}^k \begin{aligned}[t] \{
        \langle p''', \lambda(p_h, \sigma_{i_h})^{-1} uvw
        \rangle \mid
         & \exists \langle p', u \rangle \in p_h^r                      \\
         & \exists \langle p', \sigma_{i_h}, v, p'' \rangle \in
        \Delta                                                          \\
         & \exists \langle p''', w \rangle \in \Delta_\epsilon(p'') \}.
      \end{aligned}.
    \end{gather*}
    Substituting the left-hand side we derive
    \begin{gather*}
      \langle q_{i_{h+1}}, \Delta_\wedge(\sigma_{i_1} \sigma_{i_2} \ldots
      \sigma_{i_h})^{-1} w_1 w_2 \ldots w_{i_{h+1}-1} \rangle \in\\
      \bigcup_{r=1}^k \begin{aligned}[t] \{
        \langle p''', \lambda(p_h, \sigma_{i_h})^{-1} uvw
        \rangle \mid
         & \exists \langle p', u \rangle \in p_h^r                      \\
         & \exists \langle p', \sigma_{i_h}, v, p'' \rangle \in
        \Delta                                                          \\
         & \exists \langle p''', w \rangle \in \Delta_\epsilon(p'') \}.
      \end{aligned}.
    \end{gather*}
    Factorizing
    \begin{gather*}
      \langle q_{i_{h+1}}, \Delta_\wedge(\sigma_{i_1} \sigma_{i_2} \ldots
      \sigma_{i_h})^{-1} w_1 w_2 \ldots w_{i_{h+1}-1} \rangle \in\\
      \bigcup_{j=1}^k \bigcup_{r=1}^k \begin{aligned}[t] \{
        \langle p''', \lambda(p_h^r, \sigma_{i_h})^{-1} uvw
        \rangle \mid
         & \exists \langle p', u \rangle \in p_h^r                  \\
         & \exists \langle p', \sigma_{i_h}, v, p'' \rangle \in
        \Delta                                                      \\
         & \exists \langle p''', w \rangle \in \Delta_\epsilon(p'')
        (p''' \in P_j) \}.
      \end{aligned}.
    \end{gather*}
    Finally
    \begin{gather*}
      \langle q_{i_{h+1}}, \Delta_\wedge(\sigma_{i_1} \sigma_{i_2} \ldots
      \sigma_{i_h})^{-1} w_1 w_2 \ldots w_{i_{h+1}-1} \rangle \in
      \bigcup_{j=1}^k p_{h+1}^j.
    \end{gather*}
    Otherwise, $h = l$ and if $n - i_l > 2$ then $\sigma_j = \epsilon$ for $j
      \in \interval{i_l+1}{n-1}$. Therefore,
    \begin{align*}
      \langle q_n, w_{i_l} w_{i_l+1} \ldots w_{n-1} \rangle \in \{
      \langle p''', vw \rangle \mid
       & \exists \langle q_{i_l}, \sigma_{i_l}, v, p'' \rangle \in
      \Delta                                                          \\
       & \exists \langle p''', w \rangle \in \Delta_\epsilon(p'') \}.
    \end{align*}
    By the inductive hypothesis we have
    \begin{gather*}
      \langle q_n, \Delta_\wedge(\sigma_{i_1}
      \sigma_{i_2} \ldots \sigma_{i_{l-1}})^{-1} w_1 w_2 \ldots
      w_{n-1} \rangle \in\\
      \bigcup_{r=1}^k \begin{aligned}[t] \{
        \langle p''', uvw \rangle \mid
         & \exists \langle p', u \rangle \in p_l^r                      \\
         & \exists \langle p', \sigma_{i_l}, v, p'' \rangle \in \Delta  \\
         & \exists \langle p''', w \rangle \in \Delta_\epsilon(p'') \}.
      \end{aligned}
    \end{gather*}
    Substituting the left-hand side further we get
    \begin{gather*}
      \langle q_n, \lambda(p_1, \sigma_{i_1} \sigma_{i_2} \ldots
      \sigma_{i_{l-1}})^{-1} w_1 w_2 \ldots w_{n-1} \rangle \in\\
      \bigcup_{r=1}^k \begin{aligned}[t] \{
        \langle p''', uvw \rangle \mid
         & \exists \langle p', u \rangle \in p_l^r                      \\
         & \exists \langle p', \sigma_{i_l}, v, p'' \rangle \in
        \Delta                                                          \\
         & \exists \langle p''', w \rangle \in \Delta_\epsilon(p'') \}.
      \end{aligned}
    \end{gather*}
    Removing the common output prefix we get
    \begin{gather*}
      \langle q_n, \lambda(p_l, \sigma_{i_l})^{-1} \lambda(p_1,
      \sigma_{i_1} \sigma_{i_2} \ldots \sigma_{i_{l-1}})^{-1} w_1 w_2
      \ldots w_{n-1} \rangle \in\\
      \bigcup_{r=1}^k \begin{aligned}[t] \{
        \langle p''', \lambda(p_l, \sigma_{i_l})^{-1} uvw
        \rangle \mid
         & \exists \langle p', u \rangle \in p_l^r                      \\
         & \exists \langle p', \sigma_{i_l}, v, p'' \rangle \in
        \Delta                                                          \\
         & \exists \langle p''', w \rangle \in \Delta_\epsilon(p'') \}.
      \end{aligned}
    \end{gather*}
    By \Cref{prop:lambda_derivative} we get
    \begin{gather*}
      \langle q_n, \lambda(p_1, \sigma_{i_1} \sigma_{i_2} \ldots
      \sigma_{i_l})^{-1} w_1 w_2 \ldots w_{n-1} \rangle \in\\
      \bigcup_{r=1}^k \begin{aligned}[t] \{
        \langle p''', \lambda(p_l, \sigma_{i_l})^{-1} uvw
        \rangle \mid
         & \exists \langle p', u \rangle \in p_l^r                      \\
         & \exists \langle p', \sigma_{i_l}, v, p'' \rangle \in
        \Delta                                                          \\
         & \exists \langle p''', w \rangle \in \Delta_\epsilon(p'') \}.
      \end{aligned}.
    \end{gather*}
    Substituting the left-hand side we derive
    \begin{gather*}
      \langle q_n, \Delta_\wedge(\sigma_{i_1} \sigma_{i_2} \ldots
      \sigma_{i_l})^{-1} w_1 w_2 \ldots w_{n-1} \rangle \in\\
      \bigcup_{r=1}^k \begin{aligned}[t] \{
        \langle p''', \lambda(p_l, \sigma_{i_l})^{-1} uvw
        \rangle \mid
         & \exists \langle p', u \rangle \in p_l^r                      \\
         & \exists \langle p', \sigma_{i_l}, v, p'' \rangle \in
        \Delta                                                          \\
         & \exists \langle p''', w \rangle \in \Delta_\epsilon(p'') \}.
      \end{aligned}.
    \end{gather*}
    Factorizing
    \begin{gather*}
      \langle q_n, \Delta_\wedge(\sigma_{i_1} \sigma_{i_2} \ldots
      \sigma_{i_l})^{-1} w_1 w_2 \ldots w_{n-1} \rangle \in\\
      \bigcup_{j=1}^k \bigcup_{r=1}^k \begin{aligned}[t] \{
        \langle p''', \lambda(p_l, \sigma_{i_l})^{-1} uvw
        \rangle \mid
         & \exists \langle p', u \rangle \in p_l^r                  \\
         & \exists \langle p', \sigma_{i_l}, v, p'' \rangle \in
        \Delta                                                      \\
         & \exists \langle p''', w \rangle \in \Delta_\epsilon(p'')
        (p''' \in P_j) \}.
      \end{aligned}.
    \end{gather*}
    Finally
    \begin{gather*}
      \langle q_n, \Delta_\wedge(\sigma_{i_1} \sigma_{i_2} \ldots
      \sigma_{i_l})^{-1} w_1 w_2 \ldots w_{n-1} \rangle \in
      \bigcup_{j=1}^k p_{l+1}^j.
    \end{gather*}
    We have $\langle q_n, \Delta_\wedge(\sigma_{i_1} \sigma_{i_2} \ldots
      \sigma_{i_l})^{-1} w_1 w_2 \ldots w_{n-1} \rangle \in
      \bigcup_{j=1}^k{p_{l+1}^j}$ and $q_n \in F$. Therefore,
    $p_{l+1}^j \cap F \neq \emptyset$ and thus $p_{l+1}^j \in F_j'$ for some $j
      \in \interval{1}{k}$. Moreover,
    \begin{align*}
      \psi_j(p_{l+1}^j)
       & = \Delta_\wedge(\sigma_{i_1} \sigma_{i_2} \ldots
      \sigma_{i_l})^{-1} w_1 w_2 \ldots w_{n-1}                             \\
       & = \lambda(p_1, \sigma_{i_1} \sigma_{i_2} \ldots \sigma_{i_l})^{-1}
      w_1 w_2 \ldots w_{n-1}.
    \end{align*}
    Therefore,
    \begin{align*}
       & \langle \sigma_{i_1} \sigma_{i_2} \ldots \sigma_{i_l}, \lambda(p_1,
      \sigma_{i_1} \sigma_{i_2} \ldots \sigma_{i_l}) \lambda(p_1,
      \sigma_{i_1} \sigma_{i_2} \ldots \sigma_{i_l})^{-1}
      w_1 w_2 \ldots w_{n-1} \rangle =                                       \\
       & \langle \sigma_{i_1} \sigma_{i_2} \ldots \sigma_{i_l}, w_1 w_2
      \ldots w_{n-1} \rangle \in \mathcal{F}_\mathcal{T'}.
    \end{align*}






    Similarly, we show that every pair of words recognized by $\mathcal{T'}$
    is also recognized by $\mathcal{T}$. Let
    \begin{align*}
      \langle \sigma_1 \sigma_2 \ldots \sigma_{n-1}, w_1 w_2 \ldots
      w_{n-1} w_n \rangle \in \mathcal{F}_\mathcal{T'}
    \end{align*}
    for $n \in \mathbb{N}$. Let
    \begin{align*}
      \langle p_1, \sigma_1, w_1, p_2 \rangle, \langle p_2, \sigma_2, w_2,
      p_3 \rangle, \ldots, \langle p_{n-1}, \sigma_{n-1}, w_{n-1}, p_n
      \rangle
    \end{align*}
    be a successful path in $\mathcal{T'}$. There exists $j \in \interval{1}{k}$
    such that $p_n^j \in F_j'$ and $w_n = \psi_j(p_n^j)$. Therefore, there
    exists $\langle q_n, u_n \rangle \in p_n^j$ such that $q_n \in F$ and $u_n
      = w_n$. Let us assume there exist $j \in \interval{1}{k}$, $q_i \in Q$ and
    $u_i, v_i \in \Omega^*$ such that $\langle q_i, u_i \rangle \in p_i^j$,
    $u_i v_i = w_i u_{i+1}$ and $\langle q_i, \sigma_i, v_i, q_{i+1} \rangle
      \in \Delta^*$ for $i \in \interval{1}{n-1}$. We have
    \begin{align*}
      p_i^j & = \delta_j(p_{i-1}, \sigma_{i-1}) \\
            & = \begin{aligned}[t]
         & \begin{aligned}[t] \{
           & \langle p''', w_{i-1}^{-1}uvw \rangle \mid           \\
           & \exists \langle p', u \rangle \in p_{i-1}^j          \\
           & \exists \langle p', \sigma_{i-1}, v, p'' \rangle \in
          \Delta                                                  \\
           & \exists \langle p''', w \rangle \in
          \Delta_\epsilon(p'')(p''' \in P_j) \} \cup
        \end{aligned} \\
         & \bigcup_{r = 1}^k{
        \begin{aligned}[t] \{
           & \langle p''', w_{i-1}^{-1}uvw \rangle \mid  \\
           & \exists \langle p', u \rangle \in p_{i-1}^r \\
           & \exists \langle p', \sigma_{i-1}, v, p''
          \rangle \in \Delta                             \\
           & \exists \langle p''', w \rangle \in
          \Delta_\epsilon(p'')(p''' \in P_r) \}
        \end{aligned}
        }
      \end{aligned}      \\
            & = \bigcup_{r = 1}^k{
      \begin{aligned}[t] \{
         & \langle p''', w_{i-1}^{-1}uvw \rangle \mid  \\
         & \exists \langle p', u \rangle \in p_{i-1}^r \\
         & \exists \langle p', \sigma_{i-1}, v, p''
        \rangle \in \Delta                             \\
         & \exists \langle p''', w \rangle \in
        \Delta_\epsilon(p'')(p''' \in P_j) \}.
      \end{aligned}
      }
    \end{align*}
    Since $\langle q_i, u_i \rangle \in p_i^j$ then $|p_i^j| > 0$. Then there
    exists $r \in \interval{1}{k}$ such that
    \begin{align*}
      \langle q_i, u_i \rangle \in \{ & \langle p''',
      w_{i-1}^{-1}uvw \rangle \mid                                                                  \\
                                      & \exists \langle p', u \rangle \in p_{i-1}^r                 \\
                                      & \exists \langle p', \sigma_{i-1}, v, p'' \rangle \in \Delta \\
                                      & \exists \langle p''', w \rangle \in \Delta_\epsilon(p'')
      (p''' \in P_j) \}
    \end{align*}
    and thus $|p_{i-1}^r| > 0$. Therefore, there exist $q_{i-1} \in Q$ and
    $u_{i-1} \in \Omega^*$ such that $\langle q_{i-1}, u_{i-1} \rangle \in
      p_{i-1}^r$.
    Moreover, there exists $v_{i-1} \in \Omega^*$ such that $\langle q_{i-1},
      \sigma_{i-1}, v_{i-1}, q_i \rangle \in \Delta^*$ and $u_i =
      w_{i-1}^{-1}u_{i-1}v_{i-1}$. Thus, $u_{i0} v_{i-1} = w_{i-1}u_i$.
    We have $v_i = u_i^{-1} w_i u_{i+1}$ for $i \in \interval{1}{n-1}$ and
    $u_n = w_n$. Therefore,
    \begin{align*}
      v_1 v_2 \ldots v_{n-1}
       & = u_1^{-1} w_1 u_2 u_2^{-1} w_2 u_3 \ldots u_{n-1}^{-1} w_{n-1}
      u_n                                                                \\
       & = u_1^{-1} w_1 w_2 \ldots w_{n-1} u_n                           \\
       & = u_1^{-1} w_1 w_2 \ldots w_n.
    \end{align*}
    By definition there exists $q_0 \in S$ such that $\langle q_0, u_1 \rangle
      \in s = p_1$ and thus $\langle q_0, \epsilon, u_1, q_1 \rangle \in \Delta^*$.
    Therefore, $\langle q_0, \sigma_1 \sigma_2 \ldots \sigma_{n-1}, w_1 w_2
      \ldots w_n, q_n \rangle \in \Delta^*$ and
    \begin{align*}
      \langle \sigma_1 \sigma_2 \ldots \sigma_{n-1}, w_1 w_2 \ldots
      w_{n-1} w_n \rangle \in \mathcal{F}_\mathcal{T}.
    \end{align*}
  \end{proof}
\fi
